\cotefile{39}
\scanpage{3}
% 1
$k$ un \deleted{corps \inserted{\deleted{anneau}} anneau avec} corps

% 2–3
$S_n(k)=k[[t_1,\ldots,t_n]]$\qquad $(n\geq0\ ;\ S_0(k)=k)$\qquad pour $n$ variable, on a anneau cosimplicial

% 4
$A_n\subset S_n(k)$\qquad sous-anneau cosimplicial i.e.

% 5
\margin{$\Big\downarrow$}% curved relocation arrow from (b) to the space after the three clauses
(b) $\forall n\geq0$, $A_n$ est un sous anneau unifère de $S_n$

% 6–9
(a) On a\quad $a_1)$ $A_n$ est stable par $\unsure{\mathfrak S|\mathscr S}_n$ opérant sur $\unsure{S_n|\text{les}}$

\hspace*{2em}$a_2)$ Si $f(t_1,\ldots,t_n)\in A_n$, alors $f(\underbrace{t_1,\ldots,t_n}_{\text{\inserted{$t_{n+1}$ variable qui ne figure pas}}})\in A_{n+1}$

\hspace*{2em}$a_3)$ Si $f(t_1,\ldots,t_n)\in A_n$, alors $f(t_1,t_2,\ldots,\underbrace{t_{n-1},t_{n-1}}_{\text{répété}})\in A_{n-1}$

% 10–11
Le plus petit $A_*=(A_n)_{n\in\mathbb N}$ \deleted{\unsure{pareil|pareille|famille}} \inserted{satisfaisant a) b)} est donné par \deleted{$A_n=\unsure{\mathbb P|\mathbb F}$}
\[
 A_n=\mathrm{Im}(\mathbb Z\to S_n)
\]

% 12–13
donc l'\deleted{objet}\inserted{anneau} cosimplicial constant de valeur $\mathrm{Im}(\mathbb Z\to k)$ (qui est $\mathbb Z$ ou un $\mathbb Z/p\mathbb Z$, \deleted{\illegible{2}}\inserted{$p$ nb premier} $=\mathrm{car}.\,k$)

% 14–17
On supposera de plus que

c)\qquad Si $f\in A_n$, \inserted{et} $f\in\mathrm{rad}\, S_n$, i.e. \deleted{\unsure{est}} $f$ inv. dans $S_n$, alors $\frac1f\in S_n$

$\Updownarrow$

c')\qquad $A_n$ est un anneau local, et l'inclusion $A_n\to S_n$ est un hom local

% 18–21
Le plus petit $A_*$ satisfaisant a) b) c) est formé des
\[
 A_n=\mathrm{Im}(\mathbb P\to S_n)
\]
où $P$ est le sous-corps premier de $k$ ($\mathbb Q$ en car $0$, $\mathbb F_p$ en car $p>0$).

% 22
d)\qquad $t_1\in A_1\subset S[[t_1]]$

% 23–24
Le plus petit $A_*$ satisfaisant à a) b) d) \deleted{\unsure{est|c)}} est formé des
\[
 A_n=\mathbb P[t_1,\ldots,t_n]_{\deleted{\unsure{\mathfrak m_n}}}\qquad
 \deleted{\unsure{\mathfrak m_n=(}}
\]

% 25–26
Le plus petit $A$ satisfaisant a) b) c) d) est formé des
\[
 A_n=\mathbb P[t_1,\ldots,t_n]_{\mathfrak m_n}\qquad
 \mathfrak m_n=\sum t_i\,\mathbb P[t_1,\ldots,t_n].
\]

% 27–30
\deleted{\unsure{a)|e)}} Un hom $S_n\longrightarrow S_m$ \inserted{ou $\underset{\unsure{\text{``}D_m\text{''}|D'_m|X_m}}{\mathrm{Specf}(S_m)}-\underset{\unsure{\text{``}D_n\text{''}|D''_n|X_n}}{\mathrm{Specf}(S_n)}$} donné par $f_1,\ldots,f_n\in S_m$ est dit un $A$-hom si $f_1,\ldots,f_n\in A_m$. Il résulte de b) et d) que l'identité $D_n\to D_n$, \deleted{est un \unsure{et les}}\inserted{\deleted{\unsure{tout applicati.|toute application}}plus gén. les applications} $D_m\to D_n$ sont \deleted{\unsure{des|donc}}\inserted{des} $A$-hom ($\varepsilon_i\in A_n$ pour $1\leq i\leq n$). L'analogue de b)
% archive sheet number: 2
